%! Representing a Categorical Variable — Unit 1, Lesson 1.2 — Slides (Beamer)
\documentclass[aspectratio=169, 11pt]{beamer}
\usepackage{apstats-beamer}   % provides \navyheader{} and \sectionlabel[color]{}
\usetikzlibrary{positioning}  % -beamer loads tikz but not this library

% The C and Y column types live in apstats-article, which a beamer deck does not
% load — redeclare them here rather than pulling in the article preamble.
\newcolumntype{C}[1]{>{\centering\arraybackslash}p{#1}}
\newcolumntype{Y}{>{\centering\arraybackslash}X}

\begin{document}

% TITLE SLIDE (CourseName is not defined in beamer, so the name is literal)
{
\setbeamercolor{background canvas}{bg=navy}
\begin{frame}[plain]
  \vspace{2.8cm}\hspace{0.55cm}%
  \begin{minipage}{0.9\textwidth}
    {\color{goldacc}\bfseries\small LESSON 1.2}\\[0.5cm]
    {\color{white}\bfseries\LARGE Representing a Categorical Variable}\\[1.2cm]
    {\color{skymid}\small Statistics~$\cdot$~Unit 1}
  \end{minipage}
\end{frame}
}

% GRADUAL RELEASE deck: title → targets → warm-up hook → I DO (notes frames) → WE DO (guided
% practice) → YOU DO (independent practice) → spoken debrief → homework → AP Practice → close.
% There is NO group activity in this course — the notes carry the whole release. There is no
% transfer set and no reflection box: the recycling audit lives only in the AP Practice page.
% Vocabulary is given up front here, not withheld: the notes frames ARE the direct instruction.
\newcommand{\redink}[1]{{\color{redacc}\bfseries #1}}

% ── LEARNING TARGETS ─────────────────────────────────────────────────────────
\begin{frame}[t, plain]
  \navyheader{Today's targets}
  \sectionlabel{Lesson 1.2}
  By the end of today, I can\ldots
  \begin{itemize}\setlength{\itemsep}{5pt}
    \item build a \textbf{frequency} and a \textbf{relative frequency} table from one column.
    \item read a \textbf{bar graph} and say which number each bar height is showing.
    \item explain why groups of \textbf{different sizes} can only be compared with rates.
    \item decide whether ``most people chose this'' is actually supported.
  \end{itemize}
  \begin{block}{How today works}
    Warm-up $\to$ \textbf{notes together} $\to$ \textbf{one we work as a class} $\to$
    \textbf{three you work alone} $\to$ debrief.
  \end{block}
\end{frame}

% ── HOOK ─────────────────────────────────────────────────────────────────────
{
\setbeamercolor{background canvas}{bg=navy}
\begin{frame}[plain]
  \vspace{1.3cm}
  \begin{center}
    {\color{skymid}\bfseries\small WHICH SCHOOL WALKS MORE?}\\[0.8cm]
    {\color{white}\bfseries\Large North: 80 walkers \qquad South: 150 walkers}\\[0.9cm]
    {\color{goldacc}\small Careful. Both answers to that question are defensible,
    and they disagree.}
  \end{center}
\end{frame}
}

% ── I DO — direct instruction begins ─────────────────────────────────────────
{
\setbeamercolor{background canvas}{bg=navy}
\begin{frame}[plain]
  \vspace{1.8cm}\hspace{0.8cm}%
  \begin{minipage}{0.86\textwidth}
    {\color{goldacc}\bfseries\small GUIDED NOTES \ \textbullet\ \ I DO}\\[0.7cm]
    {\color{white}\bfseries\Large Open to your Guided Notes page. Fill each blank as we
    name it --- not before.}\\[0.9cm]
    {\color{skymid}\small Four terms go in the vocabulary box today. Write each one the
    moment we use it.}
  \end{minipage}
\end{frame}
}

\begin{frame}[t, plain]
  \navyheader{Two tables from one column}
  \sectionlabel{notes section 1}
  \vspace{2pt}
  \renewcommand{\arraystretch}{1.3}
  \begin{tabularx}{\textwidth}{@{} l Y Y @{}}
    \rowcolor{sky}
    \textbf{How they get to school} & \textbf{Frequency} & \textbf{Relative frequency} \\
    \hline
    Bus  & 180 & 0.45 \quad (45\%) \\
    Car  & 120 & 0.30 \quad (30\%) \\
    Walk &  80 & 0.20 \quad (20\%) \\
    Bike &  20 & 0.05 \quad (5\%)  \\
    \hline
    \textbf{Total} & \textbf{400} & \textbf{1.00} \quad (100\%) \\
  \end{tabularx}
  \vspace{10pt}
  \begin{block}{The two totals are different promises}
    Frequencies add to the \textbf{group size}. Relative frequencies add to \textbf{1},
    always --- because every individual lands in exactly one category.
  \end{block}
\end{frame}

\begin{frame}[t, plain]
  \navyheader{``Most'' is a promise about half}
  \sectionlabel{the second trap}
  \vspace{8pt}
  \begin{columns}[t]
    \begin{column}{0.46\textwidth}
      \begin{block}{Plurality --- largest share}
        Bus, at 45\% --- bigger than any other single category.
      \end{block}
    \end{column}
    \begin{column}{0.46\textwidth}
      \begin{block}{Majority --- more than half}
        Would need more than 50\%. Bus does not have it.
      \end{block}
    \end{column}
  \end{columns}
  \vspace{16pt}
  \begin{center}
    {\large ``Most North students take the bus'' is \redink{not supported.}}\\[8pt]
    {\small\color{slate} 55\% of North students get to school some other way.}
  \end{center}
\end{frame}

\begin{frame}[t, plain]
  \navyheader{Proportion, percent, rate}
  \sectionlabel{three names for one number}
  \vspace{10pt}
  \begin{center}
    {\Large $\dfrac{80}{400} \;=\; 0.20 \;=\; 20\% \;=\;$ 1 in 5}
  \end{center}
  \vspace{14pt}
  \begin{block}{Essential Knowledge (UNC-1.B.1)}
    Percentages, relative frequencies, and rates all provide the \textbf{same information}
    as proportions.
  \end{block}
  \vspace{6pt}
  {\small Use whichever your reader will understand fastest. On the AP exam, always attach
  the context: \emph{20\% \redink{of North students} walk}.}
\end{frame}

% ── BAR GRAPH ANATOMY ────────────────────────────────────────────────────────
\begin{frame}[t, plain]
  \navyheader{Bar graphs}
  \sectionlabel{height = count, or height = proportion}
  \vspace{2pt}
  \begin{columns}[t]
    \begin{column}{0.52\textwidth}
      \begin{center}
      \begin{tikzpicture}[scale=0.78]
        \draw[linegray, very thin] (0,1) -- (5.4,1);
        \draw[linegray, very thin] (0,2) -- (5.4,2);
        \draw[linegray, very thin] (0,3) -- (5.4,3);
        \draw[->, thick] (0,0) -- (0,4.1);
        \draw[->, thick] (0,0) -- (5.6,0);
        \foreach \y/\lab in {1/50, 2/100, 3/150} {
          \draw (-0.08,\y) -- (0.08,\y) node[left, font=\tiny, xshift=-2pt] {\lab};}
        \foreach \x/\h/\lab in {0.5/3.6/Bus, 1.8/2.4/Car, 3.1/1.6/Walk, 4.4/0.4/Bike} {
          \fill[navy] (\x,0) rectangle ++(0.8,\h);
          \node[below, font=\tiny] at (\x+0.4,-0.05) {\lab};}
      \end{tikzpicture}
      \end{center}
    \end{column}
    \begin{column}{0.44\textwidth}
      \begin{itemize}\setlength{\itemsep}{6pt}
        \item Each bar is one \textbf{category}.
        \item Its height is the \textbf{count} --- or the \textbf{proportion}, if you
              relabel the axis.
        \item \redink{Gaps between bars.} The categories are separate groups, not points
              on a number line.
        \item The order of the bars carries no meaning.
      \end{itemize}
    \end{column}
  \end{columns}
\end{frame}

% ── THE CRUX ─────────────────────────────────────────────────────────────────
\begin{frame}[t, plain]
  \navyheader{Why the district's chart was unfair}
  \sectionlabel{the error this lesson exists to kill}
  \vspace{4pt}
  \renewcommand{\arraystretch}{1.45}
  \begin{tabularx}{\textwidth}{@{} l C{2.6cm} C{2.6cm} X @{}}
    \hline
    & \textbf{Walkers} & \textbf{Surveyed} & \textbf{Walking rate} \\
    \hline
    North &  80 &   400 & \redink{20\%} \\
    South & 150 & 1,000 & 15\% \\
    \hline
  \end{tabularx}
  \vspace{12pt}
  \begin{block}{The comparison rule}
    South is 2.5$\times$ bigger, so it wins on \emph{every} raw count before anyone counts
    anything. Comparing groups of different sizes requires \redink{relative frequencies}.
  \end{block}
\end{frame}

\begin{frame}[t, plain]
  \navyheader{But counts are not always wrong}
  \sectionlabel{the counterweight --- do not skip this}
  \vspace{10pt}
  \begin{columns}[t]
    \begin{column}{0.46\textwidth}
      \begin{block}{``Which school walks \emph{more often}?''}
        A question about \textbf{share}. Use the rate: North, 20\% to 15\%.
      \end{block}
    \end{column}
    \begin{column}{0.46\textwidth}
      \begin{block}{``How many students walk in this district?''}
        A question about \textbf{how many}. Use the count: 230.
      \end{block}
    \end{column}
  \end{columns}
  \vspace{16pt}
  \begin{center}
    {\large The rule is not ``always use rates.''}\\[8pt]
    {\small\color{slate} It is: \redink{say which number you used, and why that one.}}
  \end{center}
\end{frame}

% ── WE DO ────────────────────────────────────────────────────────────────────
\begin{frame}[t, plain]
  \navyheader{Guided Practice: Two Shifts}
  \sectionlabel[goldacc]{We do $\cdot$ you hold the pen}
  A coffee shop records drink size. Morning: \textbf{240} orders. Evening: \textbf{80}.
  \vspace{4pt}
  \renewcommand{\arraystretch}{1.25}
  \begin{tabularx}{\textwidth}{@{} l Y Y Y Y @{}}
    \rowcolor{sky}
    \textbf{Shift} & \textbf{Small} & \textbf{Medium} & \textbf{Large} & \textbf{Total} \\
    \hline
    Morning & 48 & 108 & 84 & 240 \\
    Evening & 18 &  30 & 32 &  80 \\
    \hline
  \end{tabularx}
  \vspace{6pt}
  \begin{block}{Which shift sells more large drinks --- and which sells them at a higher rate?}
    Then: ``which shift is \emph{busier}?'' \redink{That one wants a count.} Knowing when a
    count is the right number matters as much as knowing when it is not.
  \end{block}
\end{frame}

% ── YOU DO ───────────────────────────────────────────────────────────────────
\begin{frame}[t, plain]
  \navyheader{Independent Practice}
  \sectionlabel{You do $\cdot$ 3 items $\cdot$ on your own}
  \begin{itemize}\setlength{\itemsep}{7pt}
    \item Drama club: 18 of 40 on honor roll. Band: 30 of 100. Which club is doing better,
          and which number says so?
    \item A spirit-week poll: Decades 100, Superheroes 90, Neon 60, out of 250. Is ``most
          students chose Decades'' supported?
    \item A student divides North's 80 walkers by 1{,}400. What did he actually compute?
  \end{itemize}
  \begin{block}{Silent and alone}
    I am circulating. \emph{Item 1} is the one I am reading over your shoulder --- if you
    answered it with a count, flag me before we move on.
  \end{block}
\end{frame}

% ── DEBRIEF ──────────────────────────────────────────────────────────────────
\begin{frame}[t, plain]
  \navyheader{Debrief: two true statements}
  \sectionlabel{8 minutes $\cdot$ pens down, papers up --- we say these aloud}
  \vspace{6pt}
  \begin{columns}[t]
    \begin{column}{0.46\textwidth}
      \begin{block}{South has more walkers}
        150 against 80. True.
      \end{block}
    \end{column}
    \begin{column}{0.46\textwidth}
      \begin{block}{North has the higher walking rate}
        20\% against 15\%. Also true.
      \end{block}
    \end{column}
  \end{columns}
  \vspace{14pt}
  \begin{block}{So which school walks more?}
    \redink{It depends on the question you are answering} --- and on a fair comparison of two
    different-sized groups, the rate is the one that answers it. Now name a question from
    today that a \redink{count} answers correctly.
  \end{block}
\end{frame}

% ── HOMEWORK (started in class, scored) ─────────────────────────────────────
\begin{frame}[t, plain]
  \navyheader{Homework --- start it now}
  \sectionlabel[goldacc]{5 items $\cdot$ scored $\cdot$ due next class}
  \begin{itemize}\setlength{\itemsep}{5pt}
    \item A club of 50 picks a service project --- and a ``most'' claim to check.
    \item Two hospitals, two flu counts, two very different sizes.
    \item Two bar graphs of the same data: what changed, what did not?
    \item Four percents that add to 95\%. What went wrong?
    \item One multiple-choice item, with a one-sentence justification.
  \end{itemize}
  \begin{block}{This one is graded}
    This is your \redink{scored} homework, on paper, due next class. Start with 2, 4, and 5
    while I am still here to answer questions --- finish the rest at home.
  \end{block}
\end{frame}

% ── AP PRACTICE & PREVIEW (assigned, unscored) ──────────────────────────────
\begin{frame}[t, plain]
  \navyheader{AP Practice \& what's next}
  \sectionlabel[goldacc]{AP Practice --- assigned, not scored}
  \begin{itemize}\setlength{\itemsep}{5pt}
    \item Four multiple-choice items on a city recycling audit.
    \item One free-response set: 9th versus 12th graders on activity period.
    \item Optional: find a news story comparing two groups by raw counts.
  \end{itemize}
  \begin{block}{Not graded --- but this is what the exam looks like}
    The AP Practice page is \textbf{not scored}. Work it for the reps, and bring what you
    get stuck on to study hall or office hours. The \redink{graded} homework is the section
    behind it.
  \end{block}
\end{frame}

% ── CLOSE ────────────────────────────────────────────────────────────────────
{
\setbeamercolor{background canvas}{bg=navy}
\begin{frame}[plain]
  \vspace{1.5cm}\hspace{0.8cm}%
  \begin{minipage}{0.86\textwidth}
    {\color{goldacc}\bfseries\small BEFORE YOU GO}\\[0.7cm]
    {\color{white}\bfseries\Large When two groups are different sizes, counts answer
    ``how many.'' Only rates answer ``how common.''}\\[0.9cm]
    {\color{skymid}\small Next: Lesson 1.3 --- the values stop repeating, so counting stops
    working. Dotplots, stemplots, and histograms.}
  \end{minipage}
\end{frame}
}

\end{document}
